%%%%%%%%%%%%%%%%%%%%%%%%%%%%%%%%%%%%%%%%%
% Short Sectioned Assignment LaTeX Template Version 1.0 (5/5/12)
% This template has been downloaded from: http://www.LaTeXTemplates.com
% Original author:  Frits Wenneker (http://www.howtotex.com)
% License: CC BY-NC-SA 3.0 (http://creativecommons.org/licenses/by-nc-sa/3.0/)
%%%%%%%%%%%%%%%%%%%%%%%%%%%%%%%%%%%%%%%%%

% \documentclass[paper=a4, fontsize=11pt]{scrartcl} % A4 paper and 11pt font size
\documentclass[11pt, a4paper]{book}
\usepackage[T1]{fontenc} % Use 8-bit encoding that has 256 glyphs
\usepackage[utf8]{inputenc}
\usepackage{fourier} % Use the Adobe Utopia font for the document - comment this line to return to the LaTeX default
\usepackage{listings} % para insertar código con formato similar al editor
\usepackage[spanish, es-tabla]{babel} % Selecciona el español para palabras introducidas automáticamente, p.ej. "septiembre" en la fecha y especifica que se use la palabra Tabla en vez de Cuadro
\usepackage{url} % ,href} %para incluir URLs e hipervínculos dentro del texto (aunque hay que instalar href)
\usepackage{graphics,graphicx, float} %para incluir imágenes y colocarlas
\usepackage[gen]{eurosym} %para incluir el símbolo del euro
\usepackage{cite} %para incluir citas del archivo <nombre>.bib
\usepackage{enumerate}
\usepackage{hyperref}
\usepackage{graphicx}
\usepackage{tabularx}
\usepackage{booktabs}

\usepackage[table,xcdraw]{xcolor}
\hypersetup{
	colorlinks=true,	% false: boxed links; true: colored links
	linkcolor=black,	% color of internal links
	urlcolor=cyan		% color of external links
}
\renewcommand{\familydefault}{\sfdefault}
\usepackage{fancyhdr} % Custom headers and footers
\pagestyle{fancyplain} % Makes all pages in the document conform to the custom headers and footers
\fancyhead[L]{} % Empty left header
\fancyhead[C]{} % Empty center header
\fancyhead[R]{Jesús Miguel Jaldo Ruiz} % My name
\fancyfoot[L]{} % Empty left footer
\fancyfoot[C]{} % Empty center footer
\fancyfoot[R]{\thepage} % Page numbering for right footer
%\renewcommand{\headrulewidth}{0pt} % Remove header underlines
\renewcommand{\footrulewidth}{0pt} % Remove footer underlines
\setlength{\headheight}{13.6pt} % Customize the height of the header

\usepackage{titlesec, blindtext, color}
\definecolor{gray75}{gray}{0.75}
\newcommand{\hsp}{\hspace{20pt}}
\titleformat{\chapter}[hang]{\Huge\bfseries}{\thechapter\hsp\textcolor{gray75}{|}\hsp}{0pt}{\Huge\bfseries}
\setcounter{secnumdepth}{4}
\usepackage[Lenny]{fncychap}
\usepackage[table]{xcolor} 
\setlength{\parskip}{2mm}

\begin{document}

	% Plantilla portada UGR
	\begin{titlepage}
 
 
\newlength{\centeroffset}
\setlength{\centeroffset}{-0.5\oddsidemargin}
\addtolength{\centeroffset}{0.5\evensidemargin}
\thispagestyle{empty}

\noindent\hspace*{\centeroffset}\begin{minipage}{\textwidth}

\centering
\includegraphics[width=0.9\textwidth]{imagenes/logo_ugr.jpg}\\[1.4cm]

\textsc{ \Large TRABAJO FIN DE GRADO\\[0.2cm]}
\textsc{INGENIERÍA INFORMÁTICA}\\[1cm]
% Upper part of the page
% 
% Title
{\Huge\bfseries Covid-19 Reports\\
}
\noindent\rule[-1ex]{\textwidth}{3pt}\\[3.5ex]
%{\large\bfseries Subtitulo del Proyecto}
\end{minipage}

\vspace{2.5cm}
\noindent\hspace*{\centeroffset}\begin{minipage}{\textwidth}
\centering

\textbf{Autor}\\ {Jesús Miguel Jaldo Ruiz}\\[2.5ex]
\textbf{Director}\\
{Juan Julián Merelo Guervós}\\[2cm]
\includegraphics[width=0.3\textwidth]{imagenes/etsiit_logo.png}\\[0.1cm]
\textsc{Escuela Técnica Superior de Ingenierías Informática y de Telecomunicación}\\
\textsc{---}\\
Granada, Octubre de 2020
\end{minipage}
%\addtolength{\textwidth}{\centeroffset}
%\vspace{\stretch{2}}
\end{titlepage}




	% Plantilla prefacio UGR
	\thispagestyle{empty}

\begin{center}
{\large\bfseries Covid-19 Reports \\ Subtítulo }\\
\end{center}
\begin{center}
Jesús Miguel Jaldo Ruiz\\
\end{center}

%\vspace{0.7cm}

\vspace{0.5cm}
\noindent{\textbf{Palabras clave}: \textit{software libre, covid-19, pandemia}
\vspace{0.7cm}

\noindent{\textbf{Resumen}\\
	

\cleardoublepage

\begin{center}
	{\large\bfseries Covid-19 Reports}\\
\end{center}
\begin{center}
Jesús Miguel Jaldo Ruiz\\
\end{center}
\vspace{0.5cm}
\noindent{\textbf{Keywords}: \textit{open source, covid-19, pandemic}, \textit{floss}
\vspace{0.7cm}

\noindent{\textbf{Abstract}\\


\cleardoublepage

\thispagestyle{empty}

\noindent\rule[-1ex]{\textwidth}{2pt}\\[4.5ex]

D. \textbf{Tutora/e(s)}, Profesor(a) del ...

\vspace{0.5cm}

\textbf{Informo:}

\vspace{0.5cm}

Que el presente trabajo, titulado \textit{\textbf{Chief}},
ha sido realizado bajo mi supervisión por \textbf{Estudiante}, y autorizo la defensa de dicho trabajo ante el tribunal
que corresponda.

\vspace{0.5cm}

Y para que conste, expiden y firman el presente informe en Granada a Noviembre de 2020.

\vspace{1cm}

\textbf{El/la director(a)/es: }

\vspace{5cm}

\noindent \textbf{(nombre completo tutor/a/es)}

\chapter*{Agradecimientos}






	% Índice de contenidos
	\newpage
	\tableofcontents

	% Índice de imágenes y tablas
	\newpage
	\listoffigures

	% Si hay suficientes se incluirá dicho índice
	\listoftables 
	\newpage

	% Introducción 
	\chapter{Introducción}

Este proyecto es software libre, y está liberado con la licencia \cite{gplv3}.

\section{Motivación}

Vivimos en tiempos extraños, desde diciembre de 2019 escuchamos hablar sobre la aparición de un brote epidemiológico de neumonía cuya causa era desconocida por aquel entonces en la ciudad china de Wuhan. Este brote tuvo lugar en el mercado de la ciudad y tras el aviso de las autoridades competentes de Wuhan a la OMS (Organización Mundial de la Salud) se empezó a investigar que podría estar provocando dicho brote. Esto hizo que se tras varias investigaciones se conociera que la causa del mismo era la aparición de un nuevo coronavirus \cite{oms-covid} de tipo zoonótico, al cual se le dio el nombre de Covid-19. Según declaraciones de la OMS, este virus presentaba un riesgo para la salud pública, bajo las regulaciones del Reglamento Sanitario Internacional \cite{reglamento-sanitario-internacional} y posteriormente este seria considerado como una pandemia.

Poco a poco este nuevo virus fue extendiéndose por el resto del planeta, empezando por el continente asiático y extendiéndose por el resto del planera.Todos y cada uno de los países afectados han aportado los datos de las diferentes incidencias que ha tenido el virus dentro de sus fronteras, aunque a pesar de ello seguimos sin conocer el alcance real de la pandemia ya que no existe un criterio común para la publicación de los datos, si no que cada país los calcula siguiendo los métodos que considera oportunos, por lo que la incidencia puede ser incluso mayor de lo mostrado por los datos, pero no vamos a centrarnos en como se calculan dichos datos.

Basándonos en esto, sabemos que existe una preocupación por parte de la población hacia el virus, ala cual le gusta estar informada. Sabemos que hoy en día es importante estar bien informado, como nos muestra Gabriela Nova en su articulo "Los beneficios de estar informado" \cite{gabriela-nova}. Por ello, la mayor preocupación de la gente es conocer información a cerca del virus, de como está evolucionando, cuantas contagios, muertes y altas se han producido, entre otra información.

Como hemos dicho, a la gente le gusta estar informada, y más ahora, por lo que la transparencia en una situación como la que se está viviendo es algo más que esencial. Esta información tiene como funcionalidad mantener a la gente tranquila, pudiendo aportarles datos de la evolución de la pandemia, pero no solo eso, si no que también sirven para hacer un análisis de cara a la implementación y toma de medidas por parte de los gobiernos de los países afectados. Sin embargo todos estos datos a pesar de ser accesibles para la población, presenta un problema, la difícil visualización de esto de una manera sencilla e intuitiva. En este punto podemos hacernos una pregunta: ¿como podemos ofrecer esa información de manera que pueda ser útil a toda la población que quiera consultar y hacer uso de ella?

\section{Alcance}

\section{Objetivos generales}

	% Descripción del problema y hasta donde se llega
	\chapter{Descripción del problema}



	% Estado del arte
	% 	1. Crítica al estado del arte
	% 	2. Propuesta
	\chapter{Estado del arte}


	
	\chapter{Planificación}

\section{Metodología utilizada}

\section{Temporización}

\section{Seguimiento del desarrollo}


	% Análisis del problema
	% 1. Análisis de requisitos
	% 2. Análisis de las soluciones
	% 3. Solucion propuesta
	% 4. Análisis de seguridad
	\chapter{Análisis del problema}
 
En este punto se tiene que ser capaz de plantear los objetivos que han de ser alcanzados a lo largo del desarollo del producto final.Para hacer esto se hará uso de una serie de historias de usuarios, un elemento básico a la hora de aplicar metodologías ágiles en un proyecto y especialmente para poder aplicar SCRUM. Las historias de usuario representarán de manera breve las características demandadas por el cliente las cuales deberán formar parte de la funcionalidad del producto, satisfaciendo sus exigencias.

El proceso por el cual se realiza la extracción de información relacionada con la funcionalidad del proyecto se debe llevar a cabo entre los miembros del equipo y el propio cliente. Al aplicar SCRUM, este proceso no solo se realizará en la fase inicial del proyecto, si no que se realizan en cada sprint del proyecto, de manera que se pueda obtener el resultado esperado en un corto espacio de tiempo y permitiendo amoldar el proyecto a lo requerido por el cliente de la forma mas eficiente posible.

Las historias de usuario, en terminos generales, siempre han de extraerse durante las reuniones con el cliente y es deseable que sean escritas por el mismo y en un lenguaje claro, sin entrar en detalles. Estas han de aportarnos la funcionalidad requeriada por el proyecto, entregando de esta forma un valor particular al cliente.

Estas han de desglosarse en tres apartados:
\begin{itemize}
	\item \textbf{Como}: representa el rol que va a utilizar el proyecto
	\item \textbf{Quiere}: representa la acción o evento que quiere que ocurra
	\item \textbf{Para}: representa la funcionalidad que se quiere cubrir.
\end{itemize}

\newpage
A su vez tambien puede usarse la estructura presentada por la web Scrum Manager\cite{scrum-manager}:
\begin{itemize}
	\item Nombre breve y descriptivo.
	\item Descripción de la funcionalidad en forma de diálogo o monólogo del usuario describiendo la funcionalidad que desea realizar.
	\item Criterio de validación y verificación que determinará para considerar terminado y aceptable por el cliente el desarrollo de la funcionalidad descrita.
\end{itemize}

Como se ha dicho antes, las historias de usuario ayudan a modelar el producto según las necesidades del cliente y mediante una reunión entre el equipo y este. El cliente aportará la idea que tiene, las necesidades que pretende cubrir y las funcionalidades que en cada momento el estima oportunas para el proyecto. A su vez el equipo encargado del desarrollo también podrán aportar su punto de vista en ciertos puntos con la finalidad de poder enriquecer el proyecto. Para finalizar el Product Owner, que actúa como la voz del cliente dentro del equipo, será el encargado de redactar las historias de usuario y de extraer las diferentes tareas resultantes de las mismas, identificandolas según el coste su coste y prioridad. Con esto lo que se consigue es definir el Product Backlog, base a la hora de aplicar SCRUM a un proyecto.

Es importante recalcar que en el Product Backlog se indicaran los diferentes sprints del proyecto y las tareas asociadas a los mismos. Esto no quiere decir que se mantenga durante todo el proyecto, pues la definición obtenida al inicio del mismo puede varias dependiendo de las necesidades del cliente. Dado que el proyecto estará dividido en una serie de sprints, durante los mismos debería realizarse una reunión entre los diferentes miembros del equipo y el cliente donde se podrán hacer adaptaciones que se consideren convenientes permitiendo cambiar o replantear los objetivos del proyecto con el fin de maximizar su utilidad.

Las historias de usuario serán definidas con la siguiente estructura:


\rowcolors{1}{gray!30}{gray!10}
\begin{table}[h]
	\begin{center}
		\begin{tabular}{| c | p{9cm} |}
			\hline
			
			Historia de Usuario &  \\ \hline
		 
		 
			ID & HUXX \\
			Nombre &  \\
			Prioridad &  \\
			Riesgo &  \\
			Descripción &  \\
			Validación &  \\ \hline
		\end{tabular}
		\caption{Modelo historia de Usuario}
		\label{tab:historia-usuario}
	\end{center}
\end{table}

\newpage
Cada campo representa lo siguiente:
\begin{itemize}
	\item \textbf{ID}: Identificar único de la historia de usuario.
	\item \textbf{Nombre}: Nombre asignado a la historia de usuario
	\item \textbf{Prioridad}: Importancia a la hora de llevar a cabo en el desarrollo, pudiendo ser alta, media o baja
	\item \textbf{Riesgo}: Importancia en relación al conjunto del proyecto, indicando así en caso de fallo el daño provocado, pudiendo ser alto, medio o bajo. .
	\item \textbf{Descripción}: Explicación de la historia de usuario, dejando clara la idea de la misma
	\item \textbf{Validación}: Condiciones que se han de cumplir para dar la histordia por finalizada.
\end{itemize}


\rowcolors{1}{gray!30}{gray!10}
\begin{table}[h]
	\begin{center}
		\begin{tabular}{| c | p{9cm} |}
			\hline
			
			Historia de Usuario &  \\ \hline
			
			
			ID & HU01 \\
			Nombre & Apariencia \\
			Prioridad & Baja \\
			Riesgo & Baja \\
			Descripción & Como usuario quiero que el producto tenga un diseño, simple, sencillo e intuitivo. \\
			Validación & \begin{itemize}
							\item Quiero acceder a la información pulsando un boton.
							\item Quiero ver todos los datos de manera clara.
						\end{itemize} \\ \hline
		\end{tabular}
		\caption{Historia de Usuario Apariencia}
		\label{tab:historia-usuario-01}
	\end{center}
\end{table}

\rowcolors{1}{gray!30}{gray!10}
\begin{table}[h]
	\begin{center}
		\begin{tabular}{| c | p{9cm} |}
			\hline
			
			Historia de Usuario &  \\ \hline
			
			
			ID & HU02 \\
			Nombre & Funcionamiento \\
			Prioridad & Baja \\
			Riesgo & Baja \\
			Descripción & Como usuario quiero que el producto pueda usarlo en todo momento, a ser posible desde un dispositivo movil. \\
			Validación & \begin{itemize}
				\item Quiero que funcione sobre todo en moviles.
				\item Quiero que tenga un acceso fácil.
			\end{itemize} \\ \hline
		\end{tabular}
		\caption{Historia de Usuario Funcionamiento}
		\label{tab:historia-usuario-02}
	\end{center}
\end{table}

\rowcolors{1}{gray!30}{gray!10}
\begin{table}[h]
	\begin{center}
		\begin{tabular}{| c | p{9cm} |}
			\hline
			
			Historia de Usuario &  \\ \hline
			
			
			ID & HU03 \\
			Nombre & Seleccionar provincia \\
			Prioridad & Baja \\
			Riesgo & Baja \\
			Descripción & Como usuario quiero poder seccionar la provincia de la que quiero consultar los datos. \\
			Validación & \begin{itemize}
				\item Quiero poder seleccionar todas las provincias.
				\item Quiero poder seleccionar toda España.
			\end{itemize} \\ \hline
		\end{tabular}
		\caption{Historia de Usuario Seleccion provincia}
		\label{tab:historia-usuario-03}
	\end{center}
\end{table}

\rowcolors{1}{gray!30}{gray!10}
\begin{table}[h]
	\begin{center}
		\begin{tabular}{| c | p{9cm} |}
			\hline
			
			Historia de Usuario &  \\ \hline
			
			
			ID & HU04 \\
			Nombre & Seleccionar datos \\
			Prioridad & Alta \\
			Riesgo & Baja \\
			Descripción & Como usuario quiero poder seccionar los datos que quiero consultar. \\
			Validación & \begin{itemize}
				\item Quiero poder ver los datos mas actuales.
				\item Quiero poder ver los datos desde el inicio de la pandemia.
				\item Quiero poder ver graficas para ver los cambios en los datos.
				\item Quiero poder ver la acumulacion de casos.
				\item Quiero poder ver los fallecimientos.
				\item Quiero poder ver las altas.
				\item Quiero poder ver los casos por provincias (solo España).
				\item Quiero poder ver las incidencias más actuales por cada 100k habitantes por provincias (solo España)
				\item Quiero poder ver las incidenciasdesde el inicio por cada 100k habitantes por provincias (solo España)
				\item Quiero poder ver como afecta el virus según la edad (solo España)
				
			\end{itemize} \\ \hline
		\end{tabular}
		\caption{Historia de Usuario Selección datos}
		\label{tab:historia-usuario-04}
	\end{center}
\end{table}

\rowcolors{1}{gray!30}{gray!10}
\begin{table}[h]
	\begin{center}
		\begin{tabular}{| c | p{9cm} |}
			\hline
			
			Historia de Usuario &  \\ \hline
			
			
			ID & HU05 \\
			Nombre & Seleccionar ayuda \\
			Prioridad & Alta \\
			Riesgo & Baja \\
			Descripción & Como usuario quiero poder usar una opción para . \\
			Validación & \begin{itemize}
				\item Quiero poder ver que es cada uno de los apartados del proyecto.
			\end{itemize} \\ \hline
		\end{tabular}
		\caption{Historia de Usuario Selección ayuda}
		\label{tab:historia-usuario-05}
	\end{center}
\end{table}



	% Desarrollo bajo sprints: 
	% 	1. Permitir registros y login de usuarios
	% 	2. Desarrollo del sistema de incidencias
	% 	3. Desarrollo del sistema de denuncias administrativas y accidentes
	% 	4. Desarrollo del sistema de croquis
	%   5. Instalación de la aplicación de manera automática
	\chapter{Implementación}

La implementación del software se ha dividido en hitos. Estos, han sido definidos en Github
y cada uno de ellos contiene un grupo de \textit{issues} que se corresponden con las distintas
mejoras que se han ido incorporando al software a lo largo de su desarrollo.\\



	% Presupuesto

	% Conclusiones
	\chapter{Conclusiones}



	% Trabajos futuros


	
	\newpage
	\bibliography{bibliografia}
	\bibliographystyle{plain}
	
	
\end{document}

